\chapterpage{Introduction}
\mychapter{Introduction}

\section{Background and Motivation}

Public transportation is one of the essential components of today's urban fabric. Buses are the fastest growing mode of transport and carry the greatest proportion of daily commuters globally, especially in the developing world, which is undergoing rapid urbanisation \cite{itf_outlook_2023}. It is crucial that passenger occupancy information is accurate in order to optimize scheduling, fleet management, prevent overcrowding and ensure safety for passengers.

Typical APC systems use the hardware sensors including an infrared beam, a pressure mat or a ticketing system. The above mentioned methods are common, but they can also have drawbacks such as high installation and maintenance expenses, susceptibility to environmental conditions, and decreased accuracy with multiple passengers simultaneously entering or exiting the system \cite{mccarthy2025video}.

With some recent progress in computer vision and deep learning, camera based passenger counting systems have been developed. Many buses are already using surveillance cameras for security reasons making vision an affordable and scalable solution \cite{McCarthy2023DeepLearning}. Neural networks can accurately assess occupancy in complex visual scenes; deep learning models that integrate object detection and density estimation have presented state of the art results in crowd counting benchmarks.

But bus interior scenes have special characteristics compared to the crowd scenes. Occupants often obstruct one another in the frame, the cameras are often placed on the ceiling creating the illusion of perspective, and the passengers' heads often seem to be small in the photographs.

\section{Problem Statement}

The ability to count passengers with high accuracy is a key need in intelligent public transportation systems. Accurate passenger data allows transport authorities and bus companies to make informed decisions, optimize routes, ensure passenger safety, and monitor bus occupancy. Traditional passengers' count mechanisms like infrared sensors, pressure sensors and ticketing systems tend to be inaccurate as a result of malfunctions, simultaneous boarding and alighting, and passenger overlap.

In rececompont years, computer vision and deep learning have shown promising results in the field of crowd counting and object detection related to surveillance tasks. In the case of bus passenger counting, however, it is not so easy to apply the methods described above due to the special conditions found in the interior of a bus. The passengers are tightly packed, so occlusion is a serious issue, in which people partially or even completely cover each other. In addition, surveillance cameras are usually installed on the ceiling, resulting in perspective distortion and making passenger heads appear small, especially in the rear sections of the bus. Illuminations, movement of vehicles and frequent passenger movements add to the difficulty in accurately counting.

The existing crowd counting methods have been developed and tested for public datasets which are primarily outdoor scenes or general crowd environments, such as ShanghaiTech\cite{Zhu2020}, UCF-QNRF\cite{ucf_qnrf_dataset} and NWPU-Crowd\cite{wang2020nwpu}. The data sets are not representative enough of the visual properties and problems of the interior of a bus. Another drawback of many studies is the estimation of crowd density instead of the identification of the individuals who make up the crowd in video frames, which can cause double counting if used in a video based application and may lead to unstable counts.

Thus, it is important to test if head detection can accurately estimate the number of passengers in a bus environment, explore optimal head detection models to detect small heads, and evaluate if head detection combined with object tracking can enhance the stability of counting and avoid multiple counting of passengers among different frames.

\section{Research Objectives}

\subsection{General Objective}

The key aim of this research is to design and test a vision based passenger count system within a bus for the interior space using the head detection and object tracking method.

\subsection{Specific Objectives}

Specific objectives are given below:

\begin{itemize}

    \item To investigate whether the method of detecting people's heads is a valid method for counting the number of people in the interior of a bus.
    
    \item To study the different object detection models and select the best model to detect mini passenger head in the real bus scenes.
    
    \item To explore the possibility of density estimation based counting as an alternative to head detection in highly occluded and high-passenger-density environments where individual heads can be challenging to be accurately localized.
    
    \item To develop a density decision method that can differentiate between those bus interior scenes that contain only a few people and those that have more people.
    
    \item To create an adaptive hybrid counting model which can combine the outputs of head detection and density estimation results into a more reliable passenger count.
    
    \item To assess the accuracy of the proposed system to compare with the ground truth results obtained by manual counting in terms of standard quantitative metrics, such as Mean Absolute Error (MAE), Root Mean Squared Error (RMSE), precision and recall.
    
    % \item To evaluate the generalizability of the models trained by testing them on new, previously unseen, images and videos of the interior of the buses.

\end{itemize}


\section{Scope and Limitations}

\subsection{Scope}

In this study, a vision based passenger counting system by means of deep learning based head detection for the interior of bus is developed and evaluated. The system to be proposed will use the images from the cameras placed in the interior of the bus to estimate the number of passengers inside the bus in real time.

The following are some of the scope of this study:

\begin{itemize}[itemsep=0pt]

\item Conduct research into a passenger counting system using head detection instead of whole-body detection for better performance in high occupancy and heavily occluded bus environments.

\item Comparing and assessing several object detection models to detect small heads of passengers in bus surveillance videos.

% \item Studying the performance of multi object tracking algorithms when keeping track of the identity of passengers across multiple video frames and mitigating the double counting problem.

% Evaluating the proposed system with the commonly used metrics such as Precision, Recall, mAP, MOTA, MAE and counting accuracy.

\item Simultaneously running experiments with videos of the bus interior captured from CVC installed in the ceiling, for various lighting conditions and passenger loads.

\end{itemize}

This study only provides for the number of passengers in buses and excludes other public transport like trains, airport and subway.

\subsection{Limitations}

The proposed solution is designed to offer an accurate and reliable passenger counting solution, but there are some limitations.

\begin{itemize}[itemsep=0pt]

\item  The system is based on video that is recorded from surveillance cameras mounted on a ceiling. Detection and tracking performance may be impacted by the camera position and/or the camera's point of view.

\item Severe occlusion is a difficult problem. If there is no visibility of the head of a passenger, the head can not be detected and undercounting may occur.

% \item The effectiveness of the tracking algorithm is strongly related to the effectiveness of the object detector. Identity switches or fragmented trajectories can result from missed detections or false detections.

\item The accuracy of the detection may be negatively affected by variations of light, shadows, movement of the bus and image quality.

\item The proposed system is based on the premise that the surveillance camera cannot be moved. In this study, camera movement or vibration, beyond normal bus operations, is not taken into consideration.

\item The study is limited to passenger counting only and does not include other aspects like passenger identification, behaviour analysis, age estimation or demographic analysis.

\end{itemize}

Although the limitations of this study exist, it is hoped that the research will yield useful information on the feasibility of head detection and object tracking for accurate passenger counting within a bus environment and can be used as a basis for further intelligent transportation studies.

\section{Significance of the Study}

The study is of relevance to both academia and practice. In the academic area, the research provides a contribution to the increasing research in the field of vision based passenger counting, by examining the integration of detection and tracking of heads and objects in a bus environment. In contrast to traditional crowd counting methods, which are based on density estimation, this paper aims to detect and track each individual passenger, thereby solving problems like occlusion, small object detection and crowd counting stability.

In practice, the passenger counting system could be used to give useful information to intelligent transportation systems. The real time passenger occupancy data can help transport authorities optimize bus schedules, minimize overcrowding, enhance the quality of service and passenger safety. Additionally, many of the buses already have surveillance cameras, meaning that the proposed method could be implemented without needing to install new hardware making it a cost effective solution for public transportation systems.

The results of this study can also be used to support the development of smart city applications in which real-time monitoring and decision making based on data are crucial.

